% Fragmento público generado desde propietarios íntegros.
% Residencia: 52.13.2.
% Fuente: ../incoming/radion_monografia_20260827/manuscrito/sections/13_sintesis.tex:135-184
\subsubsection{Sistema de teoremas componentes del radión}

\begin{theorem}[Genealogía]
El objeto \(\mathfrak R_\sigma\) desciende de
\(\APP\to\TRIT\to\TPK\to X_{\mathrm{enr}}\) y no recibe el blanco metrológico
como selector.
\end{theorem}

\begin{theorem}[Bisagra]
En la carta APP de cien marcas, la costura crítica y la fase dodecafásica
segunda satisfacen \(j_{100}(1/2)=\theta_2=50^\circ\), con \(m=2\) único.
\end{theorem}

\begin{theorem}[Densidad]
Las dos hojas del sector activo \(N_6=90\), normalizadas por \(3^6=729\),
producen \(\rho_{\mathrm{tw}}=20/81\).
\end{theorem}

\begin{theorem}[Acarreo]
La resta APP de las secciones \(S_T\) y \(S_E\) converge a
\(\epsone\) con tasa de intervalo \(3^{-54}\) por retorno y cierra exactamente
la unidad.
\end{theorem}

\begin{theorem}[Acción]
La elipse positiva satisface \(\mathcal A(\theta)=\hret\theta\); un radión
barre \(\hret\) y una vuelta encierra \(\hfull\).
\end{theorem}

\begin{theorem}[Borde--interior]
La misma frontera posee acción finita \(\hfull\), mientras el interior de
Klein tiene área hiperbólica infinita.
\end{theorem}

\begin{theorem}[Forma]
El cociente \(e^{-\eta}\) se conserva a través de semiejes espectrales,
semiejes de acción, incertidumbres, impedancia y módulo del toro.
\end{theorem}

\begin{theorem}[Polaridad]
La inversión \(C^*\mapsto-C^*\) intercambia constitutivas, invierte
impedancia y conserva la velocidad normalizada.
\end{theorem}

\begin{theorem}[Modularidad]
La inversión bidireccional actúa como \(S:\tau\mapsto-1/\tau\) y conserva
\(j(\tau)\), mientras el etiquetado focal conserva la orientación que
\(j\) olvida.
\end{theorem}

% Fuente: ../incoming/radion_monografia_20260827/manuscrito/sections/A_demostraciones.tex:1-179
\subsubsection{Demostraciones algebraicas y geométricas reunidas}

\paragraph{Demostración lineal del cierre unitario}

Para facilitar una auditoría independiente, reunimos el argumento sin
interrupción narrativa.

\begin{enumerate}[label=\textbf{Paso \arabic*.},leftmargin=2.2cm]
\item La conexión nonádica genera correlacionadamente
\(\pih,e_{\HMT},\varphi_{\HMT},\alphah\).
\item La vuelta es \(T_+=2\pih\).
\item La rueda fija \(\theta_2=50^\circ\), y la carta parangular fija
\(A=1000\alphah\).
\item La sección directa es
\[
S_E=T_+\left(\frac{50}{360}+\frac{1000\alphah}{360}\right)
=2\pih\left(\frac5{36}+\frac{25}{9}\alphah\right).
\]
\item Los intervalos certificados dan \(1<S_E<2\); luego
\(D_E=S_E-1\).
\item La incidencia produce \(N_6=90\); dos hojas y el ambiente \(729\)
producen \(\rho=2N_6/729=20/81\).
\item La torsión global es
\[
\Delta_4=\frac{\pih-e_{\HMT}\log\pih}{270}
\left(1+\frac{\pih}{729}\right).
\]
\item La sección torsional es \(S_T=1+(20/81)\Delta_4\).
\item La resta APP define
\[
\epsone=S_T-S_E=\frac{20}{81}\Delta_4-D_E.
\]
\item Por sustitución,
\[
S_E-\frac{20}{81}\Delta_4+\epsone=1.
\]
\item Multiplicar por \(360/T_+=180/\pih\) produce
\[
\frac{180^\circ}{\pih}
=50^\circ+A^\circ-\frac{20}{81}\Delta_4^\circ+\epsR^\circ.
\]
\end{enumerate}

Ningún paso recibe el lado izquierdo de la última igualdad. El resultado
metrológico sólo se evalúa al final como reconocimiento.

\paragraph{Unicidad de la resta con acarreo}

Sea \(B=1000\). Para cada posición, todo entero \(u_i\) admite una división
euclídea única
\[
u_i=Bc_i+r_i,\qquad 0\le r_i<B.
\]
Por tanto,
\[
r_i=u_i\bmod B,\qquad c_i=(u_i-r_i)/B
\]
son únicos. Recorriendo las posiciones desde el extremo remoto hacia el
frente, el acarreo de salida de una posición fija la entrada de la siguiente.
Por inducción, toda la palabra de restos y carries es única.

En pseudocódigo:
\begin{verbatim}
acarreo = remote_acarreo
for i from last_block downto first_block:
    u = torsion[i] - direct[i] + acarreo
    remainder[i] = u mod 1000
    acarreo = (u - remainder[i]) / 1000
return remainder, acarreo_word
\end{verbatim}
El procedimiento no consulta el valor esperado de la diferencia.

\paragraph{Convergencia del funcional por profundidad}

Sea
\[
F(p,e)=1+\frac{20}{81}
\frac{p-e\log p}{270}\left(1+\frac{p}{729}\right)
-2p\left(\frac5{36}+\frac{25}{9}\alphah\right).
\]
En el rectángulo compacto \(3\le p\le4\), \(2\le e\le3\), \(F\) es
continuo y continuamente diferenciable. Si \(I_N^\pi\searrow\{\pih\}\) e
\(I_N^e\searrow\{e_{\HMT}\}\), la imagen intervalar
\[
F(I_N^\pi,I_N^e)
\]
forma una familia encajada cuya anchura tiende a cero. La intersección es un
único real, \(\epsone\). La monotonía de las derivadas permite evaluar los
extremos sin muestreo interior.

La tasa \(3^{-54}\) por retorno deriva del refinamiento de seis trits a lo
largo de nueve niveles. A profundidad \(45\) tríadas, la cota de anchura
\(<4.81\times10^{-130}\) certifica más de 129 cifras del valor.

\paragraph{Derivación de la métrica y del área hiperbólica de Klein}

En el disco unitario, la forma matricial de la métrica Klein es
\[
g(u,v)=\frac1{1-r^2}I
+\frac1{(1-r^2)^2}
\begin{pmatrix}u\\v\end{pmatrix}
\begin{pmatrix}u&v\end{pmatrix}.
\]
El lema del determinante para una perturbación de rango uno da
\[
\det g=(1-r^2)^{-3}.
\]
De ahí
\[
\sqrt{\det g}\,\diff u\diff v
=(1-r^2)^{-3/2}\diff u\diff v.
\]
Integrar en polares produce \eqref{eq:klein-area-R}. El límite infinito se
debe a la singularidad métrica en la frontera; la frontera sigue teniendo
área simpléctica finita porque esa medida utiliza \(\diff Q\wedge\diff P\),
no \(\sqrt{\det g}\).

\paragraph{Distancia focal en el disco de Klein}

En un diámetro, la distancia del centro a \(r\) es
\[
d_K(0,r)=\operatorname{artanh}r.
\]
Para \(r=e_f\), \eqref{eq:excentricidad} implica
\[
1-e_f^2=e^{-2\eta}.
\]
Si \(u_f=\operatorname{artanh}e_f\), entonces
\[
\operatorname{sech}u_f=\sqrt{1-e_f^2}=e^{-\eta},
\]
y por tanto \(\eta=\log\cosh u_f\), como en
\eqref{eq:focal-hyperbolic-chain}.

\paragraph{Transformación simpléctica de compresión}

Definamos
\begin{equation}
S_\eta=
\begin{pmatrix}e^{\eta/2}&0\\0&e^{-\eta/2}\end{pmatrix},
\qquad \det S_\eta=1.
\label{eq:compresión simpléctica}
\end{equation}
La estructura compleja transportada es
\begin{equation}
J_\eta=S_\eta
\begin{pmatrix}0&-1\\1&0\end{pmatrix}
S_\eta^{-1}
=\begin{pmatrix}0&-e^\eta\\e^{-\eta}&0\end{pmatrix},
\qquad J_\eta^2=-I.
\label{eq:Jeta}
\end{equation}
La curva
\[
\gamma(\theta)=\sqrt{2\hret}\,
S_\eta(\cos\theta,\sin\theta)
=e^{\theta J_\eta}\gamma(0)
\]
es exactamente la elipse de acción. El compresión simpléctica hiperbólico conserva área y
transporta la rotación compleja en vez de destruirla.

Para orientación \(\sigma\),
\[
\gamma_\sigma(\theta)
=(a_S\cos\theta,\sigma b_S\sin\theta),
\qquad
\mathcal A_\sigma(\theta)=\sigma\hret\theta.
\]
Las dos hojas comparten focos y módulo de acción; difieren en acción firmada.

\paragraph{Comprobación diofántica del peso dodecafásico de Barbero--Immirzi}

La ecuación \(4a-3b=45\) tiene soluciones enteras
\[
(a,b)=(12+3t,1+4t),\qquad t\in\Z.
\]
La condición \(|b|=1\) deja \(b=1\), pues \(b=-1\) no pertenece a la clase
\(1\pmod4\). De ahí \(a=12\). Esta comprobación reproduce el par obtenido mediante la cadena fundamental
dodecafásica: doce sectores y una única costura orientada. Los cardinales
\(90/120\), las multiplicidades de la cadena y el lector de retorno
conservan funciones distintas; el coeficiente de polo y la minimalidad
verifican aritméticamente su composición.
